\documentclass[12pt,a4paper]{article}

% ============================================================================
% PREAMBLE - COMPREHENSIVE ACADEMIC PAPER SETUP
% ============================================================================

% Page layout
\usepackage[margin=1in]{geometry}
\usepackage{setspace}
\onehalfspacing

% Typography and fonts
\usepackage[utf8]{inputenc}
\usepackage[T1]{fontenc}
\usepackage{lmodern}
\usepackage{microtype}

% Mathematics
\usepackage{amsmath}
\usepackage{amssymb}
\usepackage{amsthm}
\usepackage{mathtools}

% Tables and figures
\usepackage{booktabs}
\usepackage{longtable}
\usepackage{tabularx}
\usepackage{array}
\usepackage{multirow}
\usepackage{graphicx}
\usepackage{float}
\usepackage{caption}
\usepackage{subcaption}
\usepackage{threeparttable}

% Citations and references
\usepackage{natbib}

% Hyperlinks
\usepackage[colorlinks=true,linkcolor=blue,citecolor=blue,urlcolor=blue]{hyperref}

% Other utilities
\usepackage{pdflscape}
\usepackage{rotating}
\usepackage{enumitem}
\usepackage{xcolor}

% Theorem environments
\newtheorem{hypothesis}{Hypothesis}
\newtheorem{proposition}{Proposition}
\newtheorem{assumption}{Assumption}
\newtheorem{remark}{Remark}

% Custom commands
\newcommand{\E}{\mathbb{E}}
\newcommand{\Var}{\text{Var}}
\newcommand{\Cov}{\text{Cov}}
\newcommand{\plim}{\text{plim}}
\newcommand{\efw}{\text{EFW}}

% For notes in tables
\newcommand{\sym}[1]{\ifmmode^{#1}\else\(^{#1}\)\fi}

% ============================================================================
% DOCUMENT BEGIN
% ============================================================================

\begin{document}

% ============================================================================
% TITLE PAGE
% ============================================================================

\title{%
    Close Elections, Market Governments, and Economic Outcomes:\\
    A Regression Discontinuity Design Using Parliamentary Elections%
    \thanks{We thank participants at [university seminars] and [conference] for 
    helpful comments. [Names] provided excellent research assistance. 
    All errors are our own. Replication materials are available at [repository URL].
    This version: \today.}
}

\author{%
    Gabriel Saco\thanks{[Affiliation]. Email: [email]. Corresponding author.}
}

\date{January 2026}

\maketitle

% ============================================================================
% ABSTRACT
% ============================================================================

\begin{abstract}
\noindent
This paper provides quasi-experimental evidence on the causal effect of 
market-oriented electoral victories on economic institutions and macroeconomic 
performance. We exploit close parliamentary elections from the ParlGov database 
as a regression discontinuity (RD) design, where the running variable is the 
vote-share margin between market-oriented and non-market political blocs. 
Our identification strategy relies on the quasi-random assignment of electoral 
outcomes when margins are arbitrarily small.

We conduct a comprehensive specification search across 1,074 designs varying 
running variables, bandwidths, sample definitions, and control sets. Only four 
specifications pass all diagnostic criteria: density continuity (McCrary $z = 0.52$), 
covariate balance (minimum $p = 0.165$), pre-trend placebo tests, and first-stage 
strength ($F = 10.6$--$14.0$). All passing designs share common features: they use 
the vote-share running variable, focus on the positive-shock subsample (market 
challengers versus non-market incumbents), and employ a bandwidth of 0.10 
(10 percentage points).

Within this narrow but credible identification set, we find that market-oriented 
victories cause positive shocks to Economic Freedom of the World (EFW) indices, 
particularly in Area 5 (regulation of credit, labor, and business). These 
policy shifts translate into economically meaningful effects on cumulative GDP 
per capita of 0.91--1.40 percentage points over a 3--5 year horizon. Effects are 
robust across kernel specifications (triangular and uniform) and polynomial 
orders (first and second) but do not survive expanded control sets or alternative 
running variables.

Our results establish a local causal effect in close parliamentary elections 
within ParlGov coverage. We document clear boundary conditions: the identification 
is specific to vote-share margins, positive-shock samples, and parsimonious 
controls. This narrow scope is a feature, not a bug---it reflects the price of 
credible identification in observational macro. The findings bridge the 
institutions-and-growth literature with quasi-experimental methods, providing 
the first close-election RD estimates of economic freedom effects on 
macroeconomic outcomes.

\bigskip
\noindent
\textbf{Keywords:} Economic freedom, close elections, regression discontinuity, 
political economy, economic growth, identification

\bigskip
\noindent
\textbf{JEL Classification:} D72, O43, P16, E02, C21
\end{abstract}

\thispagestyle{empty}
\newpage
\setcounter{page}{1}

% ============================================================================
% TABLE OF CONTENTS (optional for submission)
% ============================================================================

\tableofcontents
\newpage

% ============================================================================
% SECTION 1: INTRODUCTION
% ============================================================================

\section{Introduction}
\label{sec:introduction}

Does electing market-oriented governments improve economic performance? This 
question sits at the intersection of political economy and growth economics, 
yet causal evidence remains scarce. A vast cross-country literature documents 
correlations between economic freedom---measured by indices such as the 
Economic Freedom of the World (EFW)---and various indicators of prosperity 
\citep{gwartney2004economic, hall1999why, murphy2019economic}. However, 
these correlations are plagued by endogeneity: countries that adopt 
market-friendly policies differ systematically from those that do not, and 
these differences may independently affect growth.

This paper provides quasi-experimental evidence on this question by exploiting 
close parliamentary elections as a regression discontinuity (RD) design. In 
elections where the vote-share margin between market-oriented and non-market 
political blocs is arbitrarily small, the winning coalition can be considered 
as-if randomly assigned. This quasi-random variation allows us to estimate 
causal effects of electing market-oriented governments on subsequent economic 
freedom and macroeconomic outcomes.

Our empirical strategy is design-first. We pre-specify a comprehensive grid 
of 1,074 specifications varying the running variable (vote share, seat share, 
margin-based), bandwidth (0.05, 0.10, 0.15), sample definition (positive shocks, 
negative shocks, pooled), EFW operationalization (levels, changes, areas), and 
control sets (baseline, reduced, expanded). For each specification, we evaluate 
five identification diagnostics:
%
\begin{enumerate}[noitemsep]
    \item \textbf{Density continuity:} McCrary test for manipulation at the cutoff
    \item \textbf{Covariate balance:} RD estimates on predetermined covariates
    \item \textbf{Pre-trend placebo:} RD estimates on lagged outcomes
    \item \textbf{First-stage strength:} F-statistic exceeding 10
    \item \textbf{Minimum sample size:} At least 10 observations on each side
\end{enumerate}
%
Only specifications that pass \textit{all} diagnostics enter our identification set.

The results are narrow but credible. Out of 1,074 specifications, only four 
pass all diagnostic criteria. These four ``baseline passers'' share 
striking commonalities:
%
\begin{itemize}[noitemsep]
    \item \textbf{Running variable:} Vote-share margin (not seat share or margins)
    \item \textbf{Sample:} Positive-shock subsample (market challengers vs.\ 
    non-market incumbents)
    \item \textbf{Bandwidth:} 0.10 (10 percentage points on either side)
    \item \textbf{Controls:} Baseline or reduced (not expanded)
    \item \textbf{EFW variant:} Area 5 (regulation) or long pre/post shock windows
\end{itemize}
%
The diagnostic statistics for these passers are reassuring: density 
$z$-statistic of 0.52 (well below 1.96), balance minimum $p$-value of 
0.165 (above our 0.15 threshold), and first-stage F-statistics of 10.6--14.0 
(above the \citealt{stock2005testing} threshold).

Within this identification set, we find economically meaningful effects. 
Market-oriented victories cause positive shocks to EFW, particularly in 
Area 5 (regulation of credit, labor, and business). These policy shifts 
translate into cumulative GDP per capita gains of 0.91--1.40 percentage 
points over a 3--5 year horizon. The effects are robust across kernel 
specifications (triangular vs.\ uniform) and polynomial orders (first vs.\ 
second), providing confidence in the local polynomial estimates.

Equally important is what does \textit{not} work. Alternative running 
variables (seat share, vote margin, top-2 margins) fail diagnostic criteria. 
Expanded control sets---which add covariates beyond our parsimonious 
baseline---fail uniformly. The ``negative'' sample (elections where 
non-market parties barely win against market incumbents) and pooled 
samples also fail. This asymmetry is substantively interesting: it suggests 
that the causal story is about market challengers defeating non-market 
incumbents, not the reverse.

These null results discipline our interpretation. We do not claim that 
economic freedom causally improves outcomes across all countries, all 
election types, or all EFW components. Our estimates are local to:
%
\begin{enumerate}[noitemsep]
    \item Close parliamentary elections in ParlGov coverage (primarily 
    European democracies)
    \item Vote-share margins within 10 percentage points of the cutoff
    \item Positive-shock subsample (market challengers)
    \item Regulatory freedom (EFW Area 5) or long pre/post shock windows
    \item Parsimonious control specifications
\end{enumerate}
%
This narrow scope is the price of credible identification in observational 
macroeconomics.

\subsection{Contributions}

This paper makes three contributions to the literature.

\paragraph{First, we provide the first close-election RD estimates of 
economic freedom effects on macroeconomic outcomes.} The institutions-and-growth 
literature has long relied on instruments such as colonial history 
\citep{acemoglu2001colonial}, legal origin \citep{laporta1999quality}, 
or geography \citep{sachs2001tropical}, all of which face exclusion 
restriction concerns. Close-election RD offers a complementary identification 
strategy with transparent assumptions. While \citet{lee2008randomized} 
pioneered this approach for U.S.\ congressional elections and 
\citet{ferreira2009votes} applied it to mayoral policy, we are the first 
to trace effects from electoral outcomes through economic freedom changes 
to GDP per capita outcomes.

\paragraph{Second, we implement a transparent specification search with 
pre-committed diagnostic thresholds.} Recent methodological work emphasizes 
the dangers of specification searching and the value of pre-registration 
\citep{christensen2018transparency, brodeur2020methods}. Our approach---grid 
search over specifications with pre-specified pass/fail criteria---guards 
against $p$-hacking while allowing the data to reveal which designs are 
identified. The fact that only 4 of 1,074 specifications pass is itself 
informative: it shows that identification is fragile and depends on 
specific design choices.

\paragraph{Third, we document the asymmetry in electoral effects and 
interpret it substantively.} The finding that only positive-shock elections 
(market challengers) pass diagnostics while negative-shock elections 
(market incumbents) fail suggests that electoral turnover---specifically, 
replacing non-market governments with market alternatives---is the key 
mechanism. This asymmetry has implications for theories of policy persistence 
and reform \citep{fernandez1991resistance, rodrik1996understanding}.

\subsection{What This Paper Does Not Claim}

To avoid over-interpretation, we state explicitly what this paper does 
\textit{not} claim:
%
\begin{itemize}[noitemsep]
    \item We do not claim that economic freedom causally affects growth 
    in all countries or contexts
    \item We do not claim that our results generalize to presidential 
    systems, developing countries, or elections with decisive margins
    \item We do not claim robustness to high-dimensional control specifications
    \item We do not claim effects in EFW areas other than Area 5 (regulation) 
    or long pre/post shock windows
    \item We do not claim symmetric effects for market incumbents facing 
    non-market challengers
\end{itemize}
%
The narrow framing is deliberate. It reflects what the evidence can and 
cannot support under rigorous identification standards.

\subsection{Roadmap}

The remainder of the paper proceeds as follows. Section \ref{sec:literature} 
reviews the related literature. Section \ref{sec:theory} presents a simple 
theoretical framework linking electoral outcomes to economic policy and 
growth. Section \ref{sec:data} describes the data sources and sample 
construction. Section \ref{sec:methodology} details the RD methodology, 
specification search, and diagnostic criteria. Section \ref{sec:results} 
presents the main results, including the identification set, first-stage, 
and IV estimates. Section \ref{sec:robustness} examines robustness across 
design choices. Section \ref{sec:heterogeneity} explores treatment effect 
heterogeneity. Section \ref{sec:discussion} discusses interpretation, 
limitations, and policy implications. Section \ref{sec:conclusion} concludes.

% ============================================================================
% SECTION 2: LITERATURE REVIEW
% ============================================================================

\section{Literature Review}
\label{sec:literature}

Our paper contributes to three literatures: the relationship between 
economic institutions and growth, the political economy of economic reform, 
and the methodological literature on close-election regression discontinuity 
designs.

\subsection{Economic Freedom and Growth}

The proposition that economic institutions---property rights, free markets, 
rule of law---promote long-run growth has deep roots in economics, dating 
at least to \citet{north1990institutions} and formalized in modern growth 
theory by \citet{acemoglu2005institutions}. The Economic Freedom of the 
World (EFW) index, published annually by the Fraser Institute since 1996 
(with quinquennial data back to 1970), provides a comprehensive measure 
aggregating 42 indicators across five areas: (1) size of government, 
(2) legal structure and property rights, (3) access to sound money, 
(4) freedom to trade internationally, and (5) regulation 
\citep{gwartney2023economic}.

A substantial empirical literature documents positive correlations between 
EFW and growth. \citet{gwartney2004economic} find that countries in the 
top quintile of economic freedom grew almost twice as fast as those in 
the bottom quintile. \citet{hall1999why} show that differences in 
``social infrastructure''---a concept closely related to economic 
freedom---explain much of the cross-country variation in output per 
worker. \citet{murphy2019economic} provide a meta-analysis of over 100 
studies and conclude that the relationship is robust across specifications, 
though effect sizes vary.

However, causal interpretation faces severe identification challenges. 
Countries that adopt market-friendly policies differ from others on 
multiple dimensions: colonial history, geography, culture, and political 
institutions, all of which may independently affect growth. Reverse 
causality is also plausible: rich countries may be more likely to adopt 
and maintain liberal economic policies because they can afford the 
short-term costs of reform.

Attempts to address endogeneity have relied on instrumental variables. 
\citet{acemoglu2001colonial} use settler mortality as an instrument for 
institutional quality, arguing that colonial powers established extractive 
institutions in high-mortality regions. \citet{laporta1999quality} use 
legal origin (common law vs.\ civil law) as an instrument for investor 
protection and government quality. \citet{hall1999why} instrument with 
distance from the equator and Western European language fraction.

These instruments have been influential but face criticism. \citet{albouy2012colonial} 
raises concerns about measurement error in settler mortality data. 
\citet{acemoglu2012colonial} responds but acknowledges data limitations. 
More fundamentally, all these instruments operate at the deep historical 
level, making it difficult to isolate the specific channel through which 
institutions affect outcomes. Legal origin, for example, may affect 
growth through mechanisms unrelated to economic freedom---financial 
development, bureaucratic capacity, or cultural attitudes toward authority.

Our close-election RD approach offers a complementary strategy. Rather 
than seeking a deep historical source of variation, we exploit variation 
that is quasi-random within the context of close parliamentary elections. 
This strategy cannot identify the effect of economic freedom across all 
countries and time periods, but it provides credible local estimates within 
its domain of applicability.

\subsection{Political Economy of Economic Reform}

A parallel literature examines the political determinants of economic 
policy. The classic partisan theory of \citet{hibbs1977political} posits 
that left-wing governments prioritize employment and equality while 
right-wing governments prioritize inflation control and property rights. 
Subsequent work has refined this picture, showing that partisan effects 
depend on institutional context \citep{garrett1998partisan}, electoral 
vulnerability \citep{mian2010political}, and the presence of policy space 
\citep{rodrik2000institutions}.

The question of whether ``politics matters'' for economic outcomes remains 
contested. \citet{potrafke2017partisan} reviews evidence from OECD panel 
studies and finds that partisan effects vary by policy domain: they are 
stronger for fiscal policy and labor market regulation, weaker for trade 
and monetary policy. \citet{persson2000political} develop theoretical 
models showing that electoral institutions (proportional vs.\ majoritarian) 
and constitutional rules (presidential vs.\ parliamentary) shape the 
feasibility of reform.

A key insight from this literature is that policy change requires political 
capital, and this capital is more abundant after decisive electoral 
victories. Close elections, by contrast, produce mandates that are 
contested and fragile. This suggests that treatment effects in close-election 
RD may differ from effects of decisive victories---a limitation we 
acknowledge but cannot address within our design.

Reform dynamics are also asymmetric: liberalizing reforms may be politically 
costly in the short run but generate benefits that create constituencies 
for their preservation \citep{fernandez1991resistance}. Conversely, 
reversals of liberalization may be rare because beneficiaries of reform 
organize to defend their gains. This asymmetry may partly explain our 
finding that positive-shock elections (market challengers) show effects 
while negative-shock elections (non-market challengers) do not.

\subsection{Close-Election Regression Discontinuity Designs}

The methodological foundation for our analysis is the close-election RD 
design introduced by \citet{lee2008randomized} to study incumbency advantage 
in U.S.\ House elections. The key insight is that in elections decided 
by infinitesimally small margins, the outcome is effectively determined 
by chance. This allows researchers to treat the election winner as-if 
randomly assigned and estimate causal effects by comparing outcomes just 
above and below the 50\% vote share threshold.

The validity of this design rests on several assumptions. First, there 
must be no precise manipulation: candidates cannot control their vote 
share with sufficient precision to game the cutoff. Second, all 
confounders must vary continuously through the cutoff: there should be 
no other policy discontinuity at exactly 50\%. Third, the running 
variable must be measured without error (or with error that is symmetric 
around the cutoff).

\citet{eggers2015validity} provide a comprehensive assessment of these 
assumptions. They examine 40,000 close U.S.\ elections and find no 
evidence of systematic sorting at the cutoff for most offices, though 
U.S.\ Senate elections show some suggestive evidence of incumbent 
advantage. \citet{caughey2011elections} raise concerns about historical 
U.S.\ House elections with bareknuckle manipulation early in U.S. history, 
but conclude that this is not a serious concern for surveys using modern 
data. \citet{cattaneo2015randomization} develop formal 
tests for manipulation based on density discontinuities at the cutoff.

Beyond American elections, close-election RD has been applied to study 
party effects on municipal policy \citep{ferreira2009votes, pettersson2008parties}, 
women's representation and public goods \citep{clots2012women, 
bhalotra2014health}, and legislator quality and policy \citep{hall2015does, 
besley2017gender}. Our application to cross-national parliamentary 
elections using ParlGov extends this methodology to the study of 
macroeconomic institutions.

A key difference from single-country applications is that we aggregate 
across elections in different countries and time periods. This raises 
concerns about pooling heterogeneous effects. We address this by 
clustering standard errors at the country level and by examining 
heterogeneity explicitly across election characteristics. The pooled 
estimate should be interpreted as an average effect across the elections 
in our sample, not as a parameter that applies uniformly to all elections.

\subsection{Connection to This Paper}

Our paper sits at the intersection of these three literatures. We share 
the substantive interest of the institutions-and-growth literature in 
estimating causal effects of economic freedom. We draw on the political 
economy literature's insights about the conditions under which reforms 
occur and persist. And we adopt the close-election RD methodology as 
our identification strategy.

The contribution is to bridge these literatures in a way that has not 
been done before. Prior work on close-election RD has focused on 
policy-specific outcomes (education spending, crime rates, municipal 
employment) rather than comprehensive economic freedom indices. Prior 
work on economic freedom and growth has relied on cross-country 
correlations or historical instruments rather than quasi-experimental 
designs. By combining the EFW index with parliamentary election data 
from ParlGov and the close-election RD framework, we provide a new 
form of evidence on a long-standing question.

% ============================================================================
% SECTION 3: THEORETICAL FRAMEWORK
% ============================================================================

\section{Theoretical Framework}
\label{sec:theory}

This section develops a simple theoretical framework linking electoral 
outcomes to economic policy and macroeconomic performance. The framework 
guides our empirical design and generates testable hypotheses.

\subsection{Setup}

Consider an economy indexed by country $i$ and time $t$. At time $t$, an 
election determines the governing coalition. Parties can be characterized 
by their ideological position $\theta \in [0, 1]$, where higher values 
indicate more market-oriented preferences. We partition parties into two 
blocs:
%
\begin{align}
    \mathcal{M} &= \{\text{parties with } \theta \geq \bar{\theta}\} 
    \quad \text{(market-oriented)} \\
    \mathcal{N} &= \{\text{parties with } \theta < \bar{\theta}\} 
    \quad \text{(non-market)}
\end{align}
%
where $\bar{\theta}$ is a threshold that we set empirically (corresponding 
to left-right position $\geq 5$ on the 0--10 scale in ParlGov).

The election outcome is summarized by the vote-share difference between blocs:
%
\begin{equation}
    R_{it} = \sum_{p \in \mathcal{M}} v_{pit} - \sum_{p \in \mathcal{N}} v_{pit},
\end{equation}
%
where $v_{pit}$ is the vote share of party $p$ in election $it$. The 
running variable $R_{it}$ is positive when market-oriented parties receive 
more votes collectively, and negative otherwise.

Define the treatment indicator as:
%
\begin{equation}
    D_{it} = \mathbf{1}[R_{it} \geq 0].
\end{equation}
%
When $D_{it} = 1$, market-oriented parties have won the popular vote and 
are more likely to form the government or lead a governing coalition.

\subsection{Policy Determination}

The governing coalition chooses economic policy $p_{it}$, which we 
interpret as the level of economic freedom. Policy is not fully 
determined by the governing coalition's ideology because of institutional 
constraints (constitutional rules, central bank independence, European 
Union requirements) and political constraints (coalition bargaining, 
legislative opposition, interest group pressure).

We model policy as:
%
\begin{equation}
    p_{it} = \gamma \cdot \theta^G_{it} + (1 - \gamma) \cdot \bar{p}_i + 
    \epsilon_{it},
    \label{eq:policy}
\end{equation}
%
where:
\begin{itemize}[noitemsep]
    \item $\theta^G_{it}$ is the ideological position of the governing coalition
    \item $\bar{p}_i$ is a country-specific baseline reflecting institutional 
    constraints
    \item $\gamma \in (0, 1)$ is the degree of policy flexibility available 
    to governments
    \item $\epsilon_{it}$ captures idiosyncratic shocks to policy implementation
\end{itemize}

The parameter $\gamma$ is key: it measures how much governments can move 
policy given constraints. If $\gamma$ is small, elections have little effect 
on policy regardless of ideology. If $\gamma$ is large, the ideology of the 
governing coalition matters substantially.

\subsection{Electoral Determination of Government Ideology}

In parliamentary systems, the governing coalition is typically formed by 
parties that can command a majority in the legislature. When the running 
variable $R_{it}$ is close to zero, the identity of the largest bloc---and 
hence the likely coalition leader---is sensitive to small vote share changes.

We model the government ideology as:
%
\begin{equation}
    \theta^G_{it} = \alpha + \tau D_{it} + f(R_{it}) + \eta_{it},
    \label{eq:first_stage}
\end{equation}
%
where:
\begin{itemize}[noitemsep]
    \item $\tau > 0$ is the discontinuous jump in government ideology at the 
    cutoff when market-oriented parties win
    \item $f(\cdot)$ is a smooth function capturing the relationship between 
    vote share and government ideology away from the cutoff
    \item $\eta_{it}$ is an error term capturing variation in coalition formation 
    conditional on the running variable
\end{itemize}

The key assumption is that in close elections where $R_{it} \approx 0$, 
the treatment $D_{it}$ is as-if randomly assigned. This assumption is 
plausible if:
\begin{enumerate}[noitemsep]
    \item Parties cannot precisely manipulate vote shares near the threshold
    \item All confounders vary continuously through the cutoff
    \item The running variable is measured without systematic error
\end{enumerate}

Under these conditions, the RD estimand $\tau$ identifies the causal effect 
of a market-oriented victory on government ideology.

\subsection{Production and Growth}

Output in country $i$ at time $t$ follows a standard neoclassical production 
function:
%
\begin{equation}
    Y_{it} = A_{it} \cdot K_{it}^{\alpha} L_{it}^{1 - \alpha},
\end{equation}
%
where $A_{it}$ is total factor productivity (TFP), $K_{it}$ is capital, 
$L_{it}$ is labor, and $\alpha \in (0, 1)$ is the capital share.

Following the institutions-and-growth literature, we assume that TFP 
depends on economic freedom:
%
\begin{equation}
    A_{it} = \bar{A}_i \cdot g(p_{it}),
\end{equation}
%
where $g(\cdot)$ is an increasing function capturing the productivity 
benefits of economic freedom. Higher economic freedom may increase TFP 
through:
\begin{itemize}[noitemsep]
    \item Improved resource allocation via market signals
    \item Reduced rent-seeking and corruption
    \item Greater incentives for entrepreneurship and innovation
    \item More efficient capital allocation
\end{itemize}

Log output per worker is:
%
\begin{equation}
    y_{it} \equiv \log \left( \frac{Y_{it}}{L_{it}} \right) = 
    \log \bar{A}_i + \log g(p_{it}) + \alpha \cdot \log k_{it},
\end{equation}
%
where $k_{it} = K_{it}/L_{it}$ is capital per worker. Changes in economic 
freedom affect output through the $\log g(p_{it})$ term.

\subsection{Testable Hypotheses}

Our theoretical framework generates several testable hypotheses.

\begin{hypothesis}[First Stage: Electoral Outcomes Affect Economic Freedom]
\label{hyp:first_stage}
Market-oriented electoral victories cause positive shifts in economic 
freedom: $\partial \E[\Delta \efw_{it}] / \partial D_{it} > 0$.
\end{hypothesis}

This hypothesis corresponds to the first stage of our IV design. If close 
market victories do not affect economic freedom, there is no channel through 
which they can affect growth.

\begin{hypothesis}[Reduced Form: Electoral Outcomes Affect Growth]
\label{hyp:reduced_form}
Market-oriented electoral victories cause higher cumulative GDP growth: 
$\partial \E[\Delta y_{it}] / \partial D_{it} > 0$.
\end{hypothesis}

This is the reduced-form effect of treatment on outcomes. A significant 
reduced-form effect is necessary (but not sufficient) for concluding that 
economic freedom affects growth.

\begin{hypothesis}[IV: Economic Freedom Affects Growth]
\label{hyp:iv}
Instrumented changes in economic freedom cause higher cumulative GDP growth. 
The IV estimate $\hat{\rho} = (\text{reduced form}) / (\text{first stage}) > 0$.
\end{hypothesis}

This is our main parameter of interest. It captures the causal effect of 
EFW shocks on GDP per capita, for the complier population (elections where 
market victories actually cause positive EFW shocks).

\begin{hypothesis}[Asymmetry: Effects Concentrate in Positive-Shock Elections]
\label{hyp:asymmetry}
Effects are larger in elections where market challengers defeat non-market 
incumbents than in elections where non-market challengers defeat market 
incumbents.
\end{hypothesis}

This hypothesis reflects the observation that liberalizing reforms may be 
politically easier to implement than reversals, and that reform beneficiaries 
organize to defend their gains. If true, we expect identification to be 
strongest in the ``positive shock'' subsample.

\begin{hypothesis}[Channel: Regulation as Primary Mechanism]
\label{hyp:regulation}
Effects concentrate in EFW Area 5 (regulation) rather than other areas 
(government size, legal system, money, trade).
\end{hypothesis}

This hypothesis arises because governments have more direct control over 
regulatory policy than over areas that depend on independent institutions 
(central banks, courts) or international agreements (trade). If true, we 
expect identification diagnostics to pass more often for Area 5 specifications.

\subsection{Interpretation of Local Effects}

The RD design identifies local effects at the cutoff. These effects 
apply to elections decided by very close margins and may not generalize 
to elections with decisive outcomes. There are at least two reasons 
why local effects may differ from average effects:

\paragraph{Mandate Effects.}
Governments that win by large margins may have more political capital 
to implement ambitious reforms. Close winners, by contrast, may pursue 
more cautious policies. This suggests that our RD estimates may be 
\textit{attenuated} relative to the average effect of electing 
market-oriented governments.

\paragraph{Marginal vs.\ Inframarginal Elections.}
Elections that are close may be close for a reason---e.g., because 
both blocs have similar policy positions. If so, the treatment effect 
of winning may be small because the policy distance between blocs is 
small. This also suggests attenuation.

We cannot distinguish these mechanisms with our data. Readers should 
interpret our estimates as applying to the population of close elections 
and exercise caution in extrapolating to the population of all elections.

% ============================================================================
% SECTION 4: DATA
% ============================================================================

\section{Data}
\label{sec:data}

Our empirical analysis combines three data sources: the ParlGov database 
for electoral and political data, the Fraser Institute's Economic Freedom 
of the World (EFW) index for economic policy measures, and the World Bank's 
World Development Indicators (WDI) for macroeconomic outcomes.

\subsection{Electoral Data: ParlGov}

The Parliament and Government Composition Database (ParlGov) provides 
comprehensive information on elections, parties, and governments in 
democratic countries \citep{doring2021parlgov}. The database covers 
37 countries with parliamentary or semi-presidential systems, primarily 
in Europe but also including Australia, Canada, Israel, Japan, New Zealand, 
and Turkey. Coverage begins as early as 1945 for some countries and extends 
to the present.

For each parliamentary election, ParlGov includes:
\begin{itemize}[noitemsep]
    \item Vote shares for all parties receiving at least 0.5\% of the vote
    \item Seat allocations in the national legislature
    \item Cabinet composition following the election
    \item Expert-coded party positions on a left-right scale (0--10)
\end{itemize}

The left-right positioning is central to our analysis. ParlGov aggregates 
expert surveys to place each party on a 0--10 scale, where 0 represents 
the extreme left and 10 represents the extreme right. We classify parties 
as ``market-oriented'' if their left-right position is $\geq 5$ and 
``non-market'' otherwise. This threshold corresponds roughly to the 
division between center-left/left parties (social democrats, greens, 
far-left) and center-right/right parties (Christian democrats, liberals, 
conservatives).

\paragraph{Running Variable Construction.}
For each election, we construct the running variable as the difference 
in vote shares between blocs:
\begin{equation}
    R_i = \sum_{p: \text{LR}_p \geq 5} v_p - \sum_{p: \text{LR}_p < 5} v_p,
\end{equation}
where $v_p$ is the vote share of party $p$ and $\text{LR}_p$ is its 
left-right position. Positive values indicate elections where 
market-oriented parties collectively received more votes.

The treatment indicator is $D_i = \mathbf{1}[R_i \geq 0]$. Note that 
we define treatment based on collective bloc vote share, not on which 
party ultimately forms the government. This ensures that the running 
variable is predetermined at the time of voting and cannot be manipulated 
after votes are counted.

\paragraph{Sample Restrictions.}
We restrict the sample to elections from 1970--2020 to ensure overlap 
with EFW data (available from 1970). We exclude elections in which no 
party-level vote share data are available, and elections in transition 
periods (e.g., the first free election after authoritarian rule). After 
these restrictions, our full sample includes 412 elections in 37 countries.

\subsection{Economic Freedom: Fraser Institute}

The Economic Freedom of the World (EFW) index has been published by the 
Fraser Institute since 1996, with data constructed retrospectively to 
1970 at five-year intervals \citep{gwartney2023economic}. Since 2000, 
the index has been published annually. The index measures the degree 
to which a country's policies and institutions support economic freedom, 
defined as the ability to engage in voluntary economic activity without 
government interference.

The EFW index aggregates 42 indicators into five areas:
\begin{enumerate}
    \item \textbf{Area 1: Size of Government.} Measures government 
    consumption, transfers and subsidies, government enterprises, and 
    top marginal tax rates.
    
    \item \textbf{Area 2: Legal System and Property Rights.} Measures 
    judicial independence, impartiality of courts, protection of property 
    rights, military interference in law and politics, and integrity of 
    the legal system.
    
    \item \textbf{Area 3: Sound Money.} Measures money supply growth, 
    inflation variability, recent inflation rates, and freedom to own 
    foreign currency accounts.
    
    \item \textbf{Area 4: Freedom to Trade Internationally.} Measures 
    tariff rates, regulatory trade barriers, black market exchange rates, 
    and controls on capital mobility.
    
    \item \textbf{Area 5: Regulation.} Measures credit market regulations, 
    labor market regulations, and business regulations.
\end{enumerate}

Each area is scored on a 0--10 scale, with higher scores indicating 
greater economic freedom. The summary index is the unweighted average 
of the five area scores.

\paragraph{EFW Shock Operationalizations.}
We consider several operationalizations of EFW shocks around elections:
\begin{itemize}[noitemsep]
    \item \textbf{Level:} $\efw_{t+h}$ for various horizons $h$
    \item \textbf{Change:} $\Delta \efw_{t+h} = \efw_{t+h} - \efw_{t}$
    \item \textbf{Pre/Post Shock:} $\overline{\efw}_{[t+1, t+h]} - 
    \overline{\efw}_{[t-h, t-1]}$
    \item \textbf{Area-Specific:} Shocks in each of the five EFW areas
\end{itemize}

Our baseline specifications use ``pre/post shock'' measures, which 
compare the average EFW in a window after the election to the average 
EFW in a window before. Specifically, we consider ``long'' windows 
(5 years pre and post for the summary index, 3 years for Area 5) and 
``short'' windows (3 years pre and post).

\subsection{Macroeconomic Outcomes: World Bank WDI}

We obtain macroeconomic outcomes from the World Bank's World Development 
Indicators (WDI). Our primary outcome is cumulative log GDP per capita, 
defined as:
\begin{equation}
    Y^{\text{cum}}_{it,h} = \sum_{s=1}^{h} \left( \log \text{GDP}_{it+s} - 
    \log \text{GDP}_{it} \right),
\end{equation}
which measures the cumulative deviation of log GDP per capita from its 
election-year level over horizon $h$. We focus on horizons of 3 and 5 years.

We also examine secondary outcomes:
\begin{itemize}[noitemsep]
    \item GDP growth (annual percentage)
    \item Gross fixed capital formation as share of GDP
    \item Government final consumption as share of GDP
    \item Inflation (CPI, annual percentage)
    \item Trade (exports + imports) as share of GDP
    \item Domestic credit to private sector as share of GDP
    \item Unemployment rate
    \item Tax revenue as share of GDP
    \item Current account balance as share of GDP
    \item Broad money (M2) as share of GDP
\end{itemize}

All monetary values are in constant 2015 U.S.\ dollars.

\subsection{Sample Construction and Summary Statistics}

We merge the three data sources by country and year. The merge requires 
non-missing data on:
\begin{enumerate}[noitemsep]
    \item Parliamentary election with vote shares and party positions
    \item EFW summary index or area score within the relevant window
    \item GDP per capita in the election year and relevant horizons
\end{enumerate}

Table~\ref{tab:summary_stats} presents summary statistics for the baseline 
analysis sample, restricted to elections within the 0.10 bandwidth 
($|R_i| \leq 0.10$).

\begin{table}[H]
\centering
\caption{Summary Statistics for Baseline Sample}
\label{tab:summary_stats}
\begin{threeparttable}
\begin{tabular}{lcccccc}
\toprule
Variable & $N$ & Mean & S.D. & Min & Median & Max \\
\midrule
\multicolumn{7}{l}{\textit{Panel A: Running Variables and Treatment}} \\
Running variable (vote share) & 60 & 0.012 & 0.058 & --0.099 & 0.008 & 0.099 \\
Market victory ($D = 1$) & 60 & 0.533 & 0.503 & 0 & 1 & 1 \\
\midrule
\multicolumn{7}{l}{\textit{Panel B: EFW Measures}} \\
EFW summary index (level) & 60 & 7.12 & 0.89 & 4.85 & 7.21 & 8.92 \\
EFW Area 5: Regulation (level) & 60 & 6.89 & 1.02 & 4.12 & 6.94 & 8.95 \\
EFW pre/post shock (summary) & 60 & 0.18 & 0.42 & --0.85 & 0.15 & 1.24 \\
EFW pre/post shock (Area 5) & 60 & 0.22 & 0.48 & --0.92 & 0.19 & 1.45 \\
\midrule
\multicolumn{7}{l}{\textit{Panel C: Macroeconomic Outcomes}} \\
Log GDP per capita & 60 & 10.21 & 0.78 & 8.45 & 10.32 & 11.52 \\
Cumulative log GDP ($h = 3$) & 60 & 0.068 & 0.052 & --0.082 & 0.065 & 0.195 \\
GDP growth (annual \%) & 60 & 2.45 & 2.12 & --5.21 & 2.52 & 8.34 \\
Investment/GDP (\%) & 60 & 22.3 & 4.8 & 12.1 & 22.0 & 35.2 \\
Government consumption/GDP (\%) & 60 & 18.5 & 3.2 & 10.2 & 18.2 & 28.4 \\
Inflation (CPI, \%) & 60 & 4.8 & 5.2 & --1.2 & 3.1 & 25.3 \\
\midrule
\multicolumn{7}{l}{\textit{Panel D: Sample Composition}} \\
Elections below cutoff ($R < 0$) & \multicolumn{6}{c}{28} \\
Elections above cutoff ($R \geq 0$) & \multicolumn{6}{c}{32} \\
Countries represented & \multicolumn{6}{c}{24} \\
Years covered & \multicolumn{6}{c}{1972--2018} \\
\bottomrule
\end{tabular}
\begin{tablenotes}
\small
\item \textit{Notes:} Summary statistics for elections within the 0.10 
bandwidth. Running variable is the vote-share difference between market 
and non-market blocs. EFW scores are on a 0--10 scale. GDP in constant 
2015 USD.
\end{tablenotes}
\end{threeparttable}
\end{table}

The baseline sample includes 60 close elections: 28 below the cutoff 
(non-market victories) and 32 above (market victories). The sample spans 
24 countries over the period 1972--2018. The average EFW score is around 
7 on the 0--10 scale, consistent with the sample consisting primarily of 
developed democracies with relatively high economic freedom.

\paragraph{Positive vs.\ Negative Shock Subsamples.}
A key distinction in our analysis is between ``positive shock'' elections 
(where market challengers face non-market incumbents) and ``negative shock'' 
elections (where non-market challengers face market incumbents). We define 
these based on the incumbent cabinet's market share:
\begin{itemize}[noitemsep]
    \item \textbf{Positive shock:} Incumbent cabinet market share $< 0.5$
    \item \textbf{Negative shock:} Incumbent cabinet market share $\geq 0.5$
\end{itemize}

In positive-shock elections, a market victory represents a shift toward 
market orientation. In negative-shock elections, a market victory 
represents maintenance of the status quo. Our identification works in 
the positive-shock subsample, which contains approximately half of the 
close elections.

% ============================================================================
% SECTION 5: EMPIRICAL METHODOLOGY
% ============================================================================

\section{Empirical Methodology}
\label{sec:methodology}

This section describes our regression discontinuity (RD) design, 
estimation approach, diagnostic tests, and comprehensive specification 
search.

\subsection{RD Design}

Our identification strategy exploits the discontinuous relationship 
between the running variable (vote-share margin) and the probability 
of forming a market-oriented government. In close elections where 
$R_i \approx 0$, the outcome---and hence the treatment---is determined 
by factors that are essentially random: weather on election day, minor 
campaign events, or statistical noise in vote counting.

\paragraph{Identifying Assumption.}
The key assumption is that all determinants of outcomes other than 
treatment vary continuously through the cutoff. Formally:
\begin{assumption}[RD Identifying Assumption]
For all confounders $X$:
\begin{equation}
    \lim_{r \uparrow 0} \E[X | R = r] = \lim_{r \downarrow 0} \E[X | R = r].
\end{equation}
\end{assumption}
Under this assumption, any discontinuity in the outcome at the cutoff 
can be attributed causally to the treatment.

\paragraph{RD Estimand.}
The estimand is the local average treatment effect (LATE) at the cutoff:
\begin{equation}
    \tau = \lim_{r \downarrow 0} \E[Y | R = r] - \lim_{r \uparrow 0} \E[Y | R = r].
    \label{eq:rd_estimand}
\end{equation}
This is the causal effect of crossing the cutoff---receiving a market 
victory rather than a non-market victory---for elections with running 
variable exactly at zero.

\subsection{Estimation}

We estimate $\tau$ using local polynomial regression. For a given 
bandwidth $h$, we restrict the sample to elections within the window 
$|R_i| \leq h$ and estimate:
\begin{equation}
    Y_i = \alpha + \tau D_i + \beta_1 R_i + \beta_2 D_i \cdot R_i + 
    X_i'\gamma + \varepsilon_i,
    \label{eq:rd_estimation}
\end{equation}
where:
\begin{itemize}[noitemsep]
    \item $Y_i$ is the outcome (EFW shock or cumulative GDP)
    \item $D_i = \mathbf{1}[R_i \geq 0]$ is the treatment indicator
    \item $R_i$ is the running variable (centered at zero)
    \item $D_i \cdot R_i$ allows different slopes on each side of the cutoff
    \item $X_i$ is a vector of baseline controls
\end{itemize}

\paragraph{Kernel Weighting.}
We weight observations using kernel functions that give more weight 
to observations close to the cutoff. We consider two kernels:
\begin{itemize}[noitemsep]
    \item \textbf{Triangular:} $K(u) = (1 - |u|) \cdot \mathbf{1}[|u| \leq 1]$
    \item \textbf{Uniform:} $K(u) = \mathbf{1}[|u| \leq 1]$
\end{itemize}
where $u = R_i / h$. The triangular kernel is asymptotically optimal 
\citep{cheng1997automatic}, while the uniform kernel provides a 
robustness check.

\paragraph{Polynomial Order.}
We primarily use local linear estimation (polynomial order 1) but also 
consider local quadratic (order 2) as a robustness check. Higher-order 
polynomials are avoided because they tend to produce erratic behavior 
at the boundary \citep{gelman2019high}.

\paragraph{Baseline Controls.}
Our baseline control set includes:
\begin{itemize}[noitemsep]
    \item Lagged log GDP per capita (election year $t$)
    \item Trade openness (election year)
    \item Inflation (election year)
    \item Lagged EFW summary index
\end{itemize}
These controls are all predetermined at the time of the election and 
serve to reduce residual variance, improving precision.

\paragraph{Standard Errors.}
We cluster standard errors at the country level to account for potential 
correlation across elections within countries. This is conservative given 
that many countries contribute only 1--3 elections to the close-election 
sample.

\subsection{Bandwidth Selection}

The choice of bandwidth involves a bias-variance tradeoff. Smaller 
bandwidths reduce bias from approximating the regression functions 
near the cutoff but increase variance due to smaller effective sample 
sizes. We address this in three ways.

\paragraph{Fixed Bandwidths.}
We consider three fixed bandwidths: $h \in \{0.05, 0.10, 0.15\}$. The 
middle value (0.10) is our primary specification; the others provide 
robustness.

\paragraph{Balance-Optimized Bandwidth.}
We implement a bandwidth selection procedure based on covariate balance. 
For each candidate bandwidth, we estimate RD effects on each predetermined 
covariate and compute the minimum $p$-value. We select the bandwidth that 
maximizes this minimum $p$-value, ensuring that no covariate is 
significantly discontinuous at the cutoff.

\paragraph{Effective Sample Size.}
We require a minimum effective sample size of 10 observations on each 
side of the cutoff. This ensures that estimates are not driven by a 
handful of observations.

\subsection{Identification Diagnostics}

We implement a comprehensive set of five diagnostic tests. A specification 
must pass \textit{all} tests to enter our identification set.

\paragraph{Test 1: Density Continuity (McCrary Test).}
If agents could precisely manipulate the running variable, we would 
expect bunching on the preferred side of the cutoff, creating a 
discontinuity in the density of $R$. We test for this using the 
\citet{mccrary2008manipulation} procedure, which estimates log-density 
on each side of the cutoff and tests whether the difference is 
statistically significant.

Our test statistic is:
\begin{equation}
    z_{\text{density}} = \frac{\log \hat{f}(0^+) - \log \hat{f}(0^-)}
    {\hat{\text{se}}}.
\end{equation}
We require $|z_{\text{density}}| < 1.96$ (i.e., no significant 
discontinuity at the 5\% level).

\paragraph{Test 2: Covariate Balance.}
If treatment is as-if randomly assigned at the cutoff, predetermined 
covariates should be balanced. We estimate equation~\eqref{eq:rd_estimation} 
with each covariate as the dependent variable and test whether the 
treatment coefficient is significant.

Let $p_k$ be the $p$-value for covariate $k$. We compute:
\begin{equation}
    p_{\min} = \min_k \{p_k\}.
\end{equation}
We require $p_{\min} > 0.15$ (i.e., no covariate is significantly 
discontinuous at even a generous threshold).

\paragraph{Test 3: Pre-Trend Placebo.}
We test whether lagged outcomes show a discontinuity. If outcomes were 
trending differentially before the election, this would suggest that 
the RD estimate captures pre-existing differences rather than treatment 
effects. We estimate RD effects on lagged outcomes (e.g., cumulative 
log GDP from $t-3$ to $t$) and require the $p$-value to exceed 0.10.

\paragraph{Test 4: First-Stage Strength.}
For our IV estimates, we require that the first-stage F-statistic 
exceeds 10, following the \citet{stock2005testing} rule of thumb. 
This guards against weak-instrument bias, which would inflate IV 
standard errors and bias point estimates toward OLS.

\paragraph{Test 5: Minimum Sample Size.}
We require at least 10 observations on each side of the cutoff:
\begin{equation}
    n_{\text{left}} \geq 10 \quad \text{and} \quad n_{\text{right}} \geq 10.
\end{equation}
This ensures that estimates are not driven by outliers.

\subsection{Two-Stage Least Squares (2SLS)}

To estimate the effect of EFW shocks on macroeconomic outcomes, we use 
electoral outcomes as an instrument for EFW changes.

\paragraph{First Stage.}
\begin{equation}
    \Delta \efw_i = \pi_0 + \pi_1 D_i + \pi_2 R_i + \pi_3 D_i \cdot R_i + 
    X_i'\delta + \nu_i.
    \label{eq:first_stage}
\end{equation}
The coefficient $\pi_1$ is the first-stage effect: the discontinuous 
jump in EFW at the cutoff.

\paragraph{Second Stage.}
\begin{equation}
    Y_i = \mu + \rho \widehat{\Delta \efw}_i + \lambda_1 R_i + 
    \lambda_2 D_i \cdot R_i + X_i'\phi + \eta_i,
    \label{eq:second_stage}
\end{equation}
where $\widehat{\Delta \efw}_i$ is the predicted value from the first 
stage. The coefficient $\rho$ is the IV estimate of the effect of EFW 
on outcomes.

\paragraph{Interpretation.}
The IV estimate $\hat{\rho}$ captures the causal effect of EFW shocks 
on outcomes, for the ``complier'' population---elections where market 
victories actually cause positive EFW shocks. If some elections are 
``defiers'' (market victories cause negative EFW shocks), this would 
violate the monotonicity assumption required for IV interpretation.

\subsection{Specification Search}

Given the many degrees of freedom in RD analysis, we implement a 
systematic specification search rather than reporting a single 
``preferred'' specification.

\paragraph{Grid of Specifications.}
We estimate results across all combinations of:
\begin{itemize}[noitemsep]
    \item \textbf{Running variable:} Vote share, seat share, vote margin, 
    top-2 margin (4 options)
    \item \textbf{Sample:} Positive shocks, negative shocks, pooled (3 options)
    \item \textbf{EFW operationalization:} Summary levels, summary changes, 
    summary pre/post, Area 5 pre/post, and other areas (10+ options)
    \item \textbf{Bandwidth:} 0.05, 0.10, 0.15, plus balance-optimized 
    (4 options)
    \item \textbf{Controls:} Baseline, reduced, expanded (3 options)
\end{itemize}

This yields a grid of 1,074 specifications for the baseline model sweep.

\paragraph{Pass/Fail Reporting.}
For each specification, we evaluate all five diagnostics and record 
whether the specification passes all tests. We report:
\begin{enumerate}[noitemsep]
    \item The number and identity of passing specifications
    \item Common features of passers (running variable, sample, bandwidth, etc.)
    \item Diagnostic statistics for each passer
    \item Effect estimates for all passers
\end{enumerate}

This approach guards against $p$-hacking by pre-committing to diagnostic 
criteria and reporting the full set of passing specifications rather 
than cherry-picking results.

\paragraph{Effect Size Calculation.}
For each passing specification, we compute the effect size as:
\begin{equation}
    \text{Effect}_{\%} = 100 \times \hat{\rho} \times \sigma_{\text{shock}},
\end{equation}
where $\hat{\rho}$ is the IV coefficient and $\sigma_{\text{shock}}$ is 
the standard deviation of the EFW shock. This converts the coefficient 
to the percentage point change in the outcome associated with a 
one-standard-deviation shock.

% ============================================================================
% SECTION 6: MAIN RESULTS
% ============================================================================

\section{Main Results}
\label{sec:results}

This section presents our main findings. We first describe the 
identification set---the specifications that pass all diagnostic 
criteria. We then present first-stage, reduced-form, and IV estimates 
for these passing specifications.

\subsection{The Identification Set}

Out of 1,074 specifications in our baseline grid search, only 
\textbf{four} pass all five diagnostic criteria. This low pass rate 
(0.37\%) reflects the stringency of our identification requirements 
and the fragility of the RD design to specification choices.

Table~\ref{tab:passers} identifies the four passing specifications 
and their key characteristics.

\begin{table}[H]
\centering
\caption{Specifications Passing All Identification Diagnostics}
\label{tab:passers}
\begin{threeparttable}
\begin{tabular}{lcccc}
\toprule
 & Spec 1 & Spec 2 & Spec 3 & Spec 4 \\
\midrule
\multicolumn{5}{l}{\textit{Panel A: Design Choices}} \\
Running variable & Vote share & Vote share & Vote share & Vote share \\
Sample & Positive & Positive & Positive & Positive \\
EFW variant & Summary & Summary & Area 5 & Area 5 \\
Pre-window & [--5, --1] & [--5, --1] & [--3, --1] & [--3, --1] \\
Post-window & [+1, +5] & [+1, +5] & [+1, +3] & [+1, +3] \\
Bandwidth & Balance-opt & Fixed 0.10 & Balance-opt & Fixed 0.10 \\
Controls & Baseline & Baseline & Baseline & Baseline \\
\midrule
\multicolumn{5}{l}{\textit{Panel B: Diagnostic Statistics}} \\
Density $z$ & 0.516 & 0.516 & 0.516 & 0.516 \\
Balance min $p$ & 0.165 & 0.165 & 0.165 & 0.165 \\
Pre-trend $p$ & $>0.10$ & $>0.10$ & $>0.10$ & $>0.10$ \\
First-stage $F$ & 10.60 & 10.60 & 14.00 & 14.00 \\
$n_{\text{left}} / n_{\text{right}}$ & 28 / 32 & 28 / 32 & 28 / 32 & 28 / 32 \\
\midrule
Pass all? & Yes & Yes & Yes & Yes \\
\bottomrule
\end{tabular}
\begin{tablenotes}
\small
\item \textit{Notes:} The four specifications that pass all identification 
diagnostics. All use vote-share running variable, positive-shock sample, 
bandwidth 0.10, and baseline controls. They differ only in EFW 
operationalization (summary vs.\ Area 5) and bandwidth selection 
(balance-optimized vs.\ fixed).
\end{tablenotes}
\end{threeparttable}
\end{table}

The four passers share striking commonalities:
\begin{itemize}[noitemsep]
    \item \textbf{Running variable:} All use vote share (not seat share, 
    not margin-based alternatives)
    \item \textbf{Sample:} All focus on positive-shock elections (market 
    challengers vs.\ non-market incumbents)
    \item \textbf{Bandwidth:} All use 0.10 bandwidth (balance-optimized 
    selects this value)
    \item \textbf{Controls:} All use baseline controls (expanded controls 
    fail uniformly)
\end{itemize}

The passers differ in EFW operationalization: two use the summary index 
with long pre/post windows ([--5, --1] and [+1, +5]), and two use Area 5 
(regulation) with shorter windows ([--3, --1] and [+1, +3]). This 
suggests that both the aggregate index and the regulation component 
show credible identification.

\paragraph{Diagnostic Interpretation.}
The density $z$-statistic of 0.516 is well below the critical value 
of 1.96, providing no evidence of sorting or manipulation at the 
cutoff. The balance minimum $p$-value of 0.165 exceeds our threshold 
of 0.15, indicating that predetermined covariates are balanced. The 
first-stage F-statistics of 10.6--14.0 exceed the Stock-Yogo threshold 
of 10, alleviating weak-instrument concerns. The sample includes 28 
elections below the cutoff and 32 above, satisfying our minimum-N 
requirement.

\subsection{What Does Not Pass}

Equally informative is what \textit{fails} the diagnostic criteria:
\begin{itemize}[noitemsep]
    \item \textbf{Alternative running variables:} Seat share, vote margin, 
    and top-2 margin specifications fail. This is concerning because it 
    suggests the identification is not robust to alternative 
    operationalizations of electoral closeness.
    
    \item \textbf{Negative-shock sample:} Elections where non-market 
    challengers face market incumbents fail diagnostics uniformly. This 
    asymmetry suggests that the causal story is about market victories 
    causing liberalization, not about non-market victories causing reversal.
    
    \item \textbf{Pooled sample:} Combining positive and negative shocks 
    fails diagnostics, likely due to heterogeneity washing out effects.
    
    \item \textbf{Expanded controls:} Adding investment, government 
    consumption, and population growth to the control set causes 
    universal failure. This may reflect over-controlling on post-treatment 
    variables or collinearity issues in small samples.
    
    \item \textbf{Alternative EFW areas:} Areas 1--4 (size of government, 
    legal system, sound money, trade) do not generate passing specifications. 
    Only Area 5 (regulation) and summary pre/post shocks pass.
\end{itemize}

These failures discipline interpretation. The identification is narrow: 
it works for a specific combination of design choices and does not 
generalize beyond this combination.

\subsection{First-Stage Results}

Table~\ref{tab:first_stage_results} presents first-stage estimates 
showing the effect of market victories on EFW shocks.

\begin{table}[H]
\centering
\caption{First-Stage Results: Effect of Market Victory on EFW Shock}
\label{tab:first_stage_results}
\begin{threeparttable}
\begin{tabular}{lcccc}
\toprule
 & (1) & (2) & (3) & (4) \\
 & Summary & Summary & Area 5 & Area 5 \\
\midrule
Market victory ($D$) & 0.342*** & 0.342*** & 0.428*** & 0.428*** \\
 & (0.105) & (0.105) & (0.114) & (0.114) \\
[0.5em]
Running variable ($R$) & 0.892 & 0.892 & 1.124 & 1.124 \\
 & (0.612) & (0.612) & (0.698) & (0.698) \\
[0.5em]
$D \times R$ & --0.445 & --0.445 & --0.621 & --0.621 \\
 & (0.824) & (0.824) & (0.912) & (0.912) \\
\midrule
Controls & Baseline & Baseline & Baseline & Baseline \\
Kernel & Triangular & Triangular & Triangular & Triangular \\
Bandwidth & 0.10 & 0.10 & 0.10 & 0.10 \\
\midrule
Observations & 60 & 60 & 60 & 60 \\
$R^2$ & 0.312 & 0.312 & 0.356 & 0.356 \\
$F$-statistic & 10.60 & 10.60 & 14.00 & 14.00 \\
\bottomrule
\end{tabular}
\begin{tablenotes}
\small
\item \textit{Notes:} First-stage regressions of EFW shocks on market-victory 
indicator. Standard errors clustered by country in parentheses. 
*** $p < 0.01$, ** $p < 0.05$, * $p < 0.10$.
\end{tablenotes}
\end{threeparttable}
\end{table}

The first-stage coefficient on market victory is 0.34 for the summary 
index and 0.43 for Area 5. These estimates indicate that market victories 
cause EFW to increase by about one-third to two-fifths of a point on 
the 0--10 scale. Both coefficients are statistically significant at the 
1\% level.

The effect on Area 5 (regulation) is larger than on the summary index, 
consistent with our hypothesis that governments have more control over 
regulatory policy than over other dimensions of economic freedom. The 
F-statistics of 10.6 and 14.0 exceed the Stock-Yogo threshold, indicating 
adequate instrument strength.

\subsection{Reduced-Form and IV Results}

Table~\ref{tab:main_results} presents our main results. Panel A shows 
reduced-form effects of market victories on cumulative log GDP per 
capita. Panel B shows IV estimates of the effect of EFW shocks on GDP.

\begin{table}[H]
\centering
\caption{Main Results: Effect on Cumulative Log GDP Per Capita}
\label{tab:main_results}
\begin{threeparttable}
\begin{tabular}{lcccc}
\toprule
 & (1) & (2) & (3) & (4) \\
 & Summary & Summary & Area 5 & Area 5 \\
\midrule
\multicolumn{5}{l}{\textit{Panel A: Reduced Form}} \\
[0.3em]
Market victory ($D$) & 0.048*** & 0.048*** & 0.039** & 0.039** \\
 & (0.018) & (0.018) & (0.016) & (0.016) \\
\midrule
\multicolumn{5}{l}{\textit{Panel B: IV Estimates}} \\
[0.3em]
EFW shock ($\widehat{\Delta \efw}$) & 0.140*** & 0.140*** & 0.091** & 0.091** \\
 & (0.045) & (0.045) & (0.038) & (0.038) \\
\midrule
Controls & Baseline & Baseline & Baseline & Baseline \\
Kernel & Triangular & Triangular & Triangular & Triangular \\
Bandwidth & 0.10 & 0.10 & 0.10 & 0.10 \\
First-stage $F$ & 10.60 & 10.60 & 14.00 & 14.00 \\
\midrule
Observations & 60 & 60 & 60 & 60 \\
Effect (\%) & 1.40 & 1.40 & 0.91 & 0.91 \\
\bottomrule
\end{tabular}
\begin{tablenotes}
\small
\item \textit{Notes:} Panel A shows reduced-form effects of market victory 
on cumulative log GDP (horizon 3 years). Panel B shows IV estimates using 
electoral outcomes as instrument for EFW shocks. Effect (\%) is 
$100 \times \hat{\rho} \times \sigma_{\text{shock}}$, the percentage 
point effect of a one-SD EFW shock on cumulative GDP. Standard errors 
clustered by country. *** $p < 0.01$, ** $p < 0.05$, * $p < 0.10$.
\end{tablenotes}
\end{threeparttable}
\end{table}

\paragraph{Reduced-Form Effects.}
The reduced-form coefficient on market victory is 0.039--0.048, indicating 
that market victories cause cumulative log GDP per capita to be 3.9--4.8 
log points higher over the post-election horizon. These effects are 
statistically significant at the 5\% level (Area 5) and 1\% level (summary).

\paragraph{IV Estimates.}
The IV estimates imply that a one-point increase in EFW causes cumulative 
log GDP to rise by 0.091--0.140 log points. To interpret magnitudes, we 
compute the effect size as $100 \times \hat{\rho} \times \sigma_{\text{shock}}$:
\begin{itemize}[noitemsep]
    \item Summary index: $100 \times 0.140 \times 0.10 = 1.40\%$
    \item Area 5: $100 \times 0.091 \times 0.10 = 0.91\%$
\end{itemize}
Thus, a one-standard-deviation EFW shock (approximately 0.1 points on 
the 0--10 scale) produces cumulative GDP gains of 0.91--1.40 percentage 
points over the 3-year horizon.

\paragraph{Economic Significance.}
These magnitudes are economically meaningful. A 1\% gain in cumulative 
GDP over 3 years translates to approximately 0.33 percentage points of 
additional annual GDP growth. Given that average annual GDP growth in 
our sample is about 2.5\%, this represents about a 13\% increase in the 
growth rate during the post-election window.

\subsection{Effects on Secondary Outcomes}

Table~\ref{tab:secondary_outcomes} presents IV estimates for secondary 
macroeconomic outcomes, using the baseline design that passes diagnostics.

\begin{table}[H]
\centering
\caption{IV Effects on Secondary Macroeconomic Outcomes}
\label{tab:secondary_outcomes}
\begin{threeparttable}
\begin{tabular}{lcc}
\toprule
Outcome & IV Coefficient & Std.\ Err. \\
\midrule
GDP growth (average \%) & 0.918** & (0.412) \\
Investment/GDP (\%) & 0.565 & (0.348) \\
Inflation (\%) & --1.338** & (0.612) \\
Government consumption/GDP (\%) & --1.838*** & (0.524) \\
Trade openness (\%) & --0.814 & (0.892) \\
Private credit/GDP (\%) & 18.11*** & (5.24) \\
Unemployment (\%) & --1.052** & (0.485) \\
Tax revenue/GDP (\%) & 1.733** & (0.712) \\
Broad money/GDP (\%) & 16.82*** & (4.89) \\
Current account (\%) & 5.225** & (2.18) \\
\bottomrule
\end{tabular}
\begin{tablenotes}
\small
\item \textit{Notes:} IV estimates of the effect of a one-point EFW 
shock on secondary outcomes. All estimates use the summary-index baseline 
design with 0.10 bandwidth and triangular kernel. Standard errors 
clustered by country. *** $p < 0.01$, ** $p < 0.05$, * $p < 0.10$.
\end{tablenotes}
\end{threeparttable}
\end{table}

The secondary outcomes show consistent patterns:
\begin{itemize}[noitemsep]
    \item \textbf{Growth and investment:} Positive EFW shocks are 
    associated with higher GDP growth and marginally higher investment
    \item \textbf{Government:} Lower government consumption and inflation, 
    consistent with market-oriented policy
    \item \textbf{Financial development:} Large increases in private 
    credit and broad money, suggesting financial liberalization
    \item \textbf{Labor market:} Lower unemployment
    \item \textbf{External:} Improved current account balance
\end{itemize}

These patterns are consistent with the expected effects of market-oriented 
reform and provide supporting evidence for the main GDP results.

% ============================================================================
% SECTION 7: ROBUSTNESS
% ============================================================================

\section{Robustness Checks}
\label{sec:robustness}

We examine the robustness of our main results to alternative design choices. 
Our specification search reveals a clear pattern: results are robust within 
the core design (varying estimation parameters) but fragile to major 
structural changes (expanded controls, alternative running variables).

\subsection{Robustness Within the Core Design}

Table~\ref{tab:robustness_estimation} shows how results vary across 
local polynomial estimation choices: kernel (triangular vs.\ uniform) 
and polynomial order (local linear vs.\ local quadratic).

\begin{table}[H]
\centering
\caption{Robustness to Kernel and Polynomial Order}
\label{tab:robustness_estimation}
\begin{threeparttable}
\begin{tabular}{lcccc}
\toprule
 & \multicolumn{2}{c}{Triangular Kernel} & \multicolumn{2}{c}{Uniform Kernel} \\
 & Order 1 & Order 2 & Order 1 & Order 2 \\
\midrule
\multicolumn{5}{l}{\textit{Panel A: EFW Summary}} \\
Effect (\%) & 1.40 & 1.40 & 1.40 & 1.40 \\
Pass diagnostics & Yes & Yes & Yes & Yes \\
\midrule
\multicolumn{5}{l}{\textit{Panel B: EFW Area 5}} \\
Effect (\%) & 0.91 & 0.91 & 0.91 & 0.91 \\
Pass diagnostics & Yes & Yes & Yes & Yes \\
\bottomrule
\end{tabular}
\begin{tablenotes}
\small
\item \textit{Notes:} Effect (\%) is the IV estimate converted to percentage 
points. All specifications use the baseline vote-share running variable, 
positive-shock sample, and 0.10 bandwidth.
\end{tablenotes}
\end{threeparttable}
\end{table}

The results are remarkably stable. The effect sizes and diagnostic pass 
rates are identical across kernel and polynomial choices for the baseline 
passers. This stability suggests that our findings are not driven by 
arbitrary choices about how to model the running variable near the cutoff.

\subsection{Fragility to Control Set Expansion}

A key limitation of our results is their sensitivity to the control set. 
While baseline controls (lagged GDP, trade, inflation, lagged EFW) yield 
passing specifications, expanding the control set causes diagnostic failure.

Table~\ref{tab:robustness_controls} summarizes pass rates by control set 
across the full grid of 1,074 specifications.

\begin{table}[H]
\centering
\caption{Robustness to Control Set Specification}
\label{tab:robustness_controls}
\begin{threeparttable}
\begin{tabular}{lcc}
\toprule
Control Set & Specs Estimated & Specs Passing \\
\midrule
Baseline & 358 & 4 \\
Reduced & 358 & 0 \\
Expanded & 358 & 0 \\
\bottomrule
\end{tabular}
\begin{tablenotes}
\small
\item \textit{Notes:} Pass rates across the full specification grid. 
Expanded controls add investment, government consumption, and population 
growth to the baseline set.
\end{tablenotes}
\end{threeparttable}
\end{table}

The failure of expanded controls likely reflects two factors:
\begin{enumerate}
    \item \textbf{Degrees of freedom:} With only ~60 observations in the 
    close-election sample, adding covariates rapidly consumes degrees of 
    freedom and reduces power.
    \item \textbf{Bad controls:} Some expanded controls (investment, 
    government consumption) may be outcomes of treatment, introducing 
    post-treatment bias.
\end{enumerate}

This fragility means we cannot claim robustness to high-dimensional 
covariate adjustment. Our identification relies on the RD assumption 
that predetermined characteristics are balanced at the limit---an assumption 
supported by our balance tests---rather than on selection-on-observables 
arguments requiring extensive controls.

\subsection{Alternative EFW Definitions}

We also tested robustness to alternative EFW definitions. As noted in 
the identification section, no specifications using Areas 1--4 pass 
diagnostics. This suggests that the ``market-oriented victory'' effect 
is not a generic shock to all dimensions of economic freedom but 
specifically a shock to regulatory policy (Area 5) and aggregate measures 
over longer horizons.

% ============================================================================
% SECTION 8: HETEROGENEITY
% ============================================================================

\section{Heterogeneity Analysis}
\label{sec:heterogeneity}

We explore heterogeneity in treatment effects by splitting the sample on 
pre-treatment characteristics. Given the small sample size (N=60 even in 
the pooled baseline), these results should be interpreted as 
\textit{exploratory} and hypothesis-generating rather than conclusive 
causal evidence.

\subsection{Heterogeneity Results}

We split the sample based on:
\begin{enumerate}
    \item \textbf{Incumbent cabinet market share:} High vs.\ Low (median split)
    \item \textbf{Ideological distance:} High vs.\ Low absolute distance 
    between blocs
\end{enumerate}

Table~\ref{tab:heterogeneity} presents the results for splits that generate 
at least one passing specification.

\begin{table}[H]
\centering
\caption{Heterogeneity Analysis (Exploratory)}
\label{tab:heterogeneity}
\begin{threeparttable}
\begin{tabular}{lcc}
\toprule
 & Low Incumbent & Low Ideological \\
 & Market Share & Distance \\
\midrule
Effect (\%) & --1.09 & --1.55 \\
First-stage $F$ & 15.54 & 14.39 \\
$n_{\text{left}} / n_{\text{right}}$ & 25 / 24 & 12 / 10 \\
Pass diagnostics? & Yes & Yes \\
\midrule
Interpretation & Sign flip & Sign flip \\
\bottomrule
\end{tabular}
\begin{tablenotes}
\small
\item \textit{Notes:} Results for subsamples that pass diagnostics. 
Effect (\%) is the standardized effect size. Note the negative sign, 
contrasting with the positive baseline effect.
\end{tablenotes}
\end{threeparttable}
\end{table}

\subsection{Interpretation of Heterogeneity}

Only two subsamples generate passing specifications: elections with low 
incumbent market share (where the incumbent is non-market) and elections 
with low ideological distance between blocs. In both cases, the estimated 
effect is \textit{negative}.

This finding is puzzling given the positive average effect in the main 
positive-shock sample. Several interpretations are possible:
\begin{itemize}
    \item \textbf{Sample selection:} The splits reduce sample size significantly 
    (e.g., to 12 vs 10 observations), making estimates volatile.
    \item \textbf{Short-term costs:} In environments with entrenched 
    non-market incumbents, market victories may trigger immediate adjustment 
    costs or political conflict that dampens growth in the 3-year window.
    \item \textbf{Policy moderation:} When ideological distance is low, 
    market victories may not represent a sharp policy discontinuity, or 
    may lead to distinct types of policy changes (e.g., fiscal consolidation) 
    that are contractionary in the short run.
\end{itemize}

Given the small samples and sign flips, we do not claim these as robust 
findings. Instead, they serve as a caution against assuming homogeneous 
effects and suggest that the positive baseline effects may be driven by 
specific electoral configurations not fully captured by simple splits.

% ============================================================================
% SECTION 9: DISCUSSION
% ============================================================================

\section{Discussion}
\label{sec:discussion}

Our results provide evidence that market-oriented electoral victories in 
close elections can have positive effects on economic freedom and 
macroeconomic performance. However, the specific pattern of results---success 
in narrow positive-shock subsamples, failure in broader samples---points 
to a nuanced interpretation.

\subsection{Interpretation of the ``Regulatory Effect''}

The most robust finding is that market victories lead to improvements 
in EFW Area 5 (regulation) and corresponding growth gains. This aligns 
with the political reality that regulatory policy is often more amenable 
to executive action than fiscal or monetary institutions.

Effect sizes of 0.91\%--1.40\% on cumulative GDP per capita over 3 years 
are substantial but plausible. They suggest that regulatory relief 
following a market victory can spur a short-to-medium-term investment 
and growth response. Note that these are \textit{cumulative} effects, 
not permanent changes to the steady-state growth rate.

\subsection{Boundary Conditions and Limitations}

We emphasize three key boundary conditions:

\paragraph{1. Local Effects.}
Our estimates apply to close elections in parliamentary democracies 
covered by ParlGov. They do not necessarily generalize to:
\begin{itemize}[noitemsep]
    \item Presidential systems (where checks and balances differ)
    \item Developing countries (where state capacity is lower)
    \item Decisive elections (where mandates are stronger)
\end{itemize}

\paragraph{2. Asymmetry.}
The failure of the ``negative shock'' sample suggests that the effect 
is not symmetric. Market victories matter when they displace non-market 
incumbents (positive shocks), but non-market victories do not appear 
to cause symmetric reversals (or at least, we cannot identify them forces). 
This is consistent with the ``ratchet effect'' in economic reform.

\paragraph{3. Fragility.}
The fragility of results to expanded controls is a genuine limitation. 
It implies that the design works only when we trust the RD assumption 
enough to rely on a parsimonious specification. Researchers should be 
cautious about pushing this identification strategy too far into 
subsample analysis or high-dimensional control spaces where data limitations 
become binding.

% ============================================================================
% SECTION 10: CONCLUSION
% ============================================================================

\section{Conclusion}
\label{sec:conclusion}

This paper presents a disciplined evaluation of the causal effect of 
electing market-oriented governments on economic outcomes. Using a 
comprehensive specification search across 1,074 designs, we identify a narrow 
but credible set of specifications where close-election RD is valid.

We find that in close parliamentary elections where market challengers 
defeat non-market incumbents, there is a causal increase in economic 
freedom (specifically regulatory policy) and a subsequent increase in 
cumulative GDP per capita of approximately 1\% over 3--5 years. These 
results are robust to estimation choices but sensitive to sample definition 
and control sets.

Our work bridges the gap between the macro-institutions literature (which 
often relies on cross-country regressions) and the micro-political economy 
literature (which uses rigorous RD designs but often for localized outcomes). 
We show that it is possible to apply credible identification to macro 
outcomes, provided one accepts the necessary boundary conditions: the effects 
are local, specific to certain institutional contexts, and cannot be 
generalized to a universal ``law'' of economic freedom.

Future research should focus on expanding the set of elections for which 
credible close-election data are available, possibly extending outside 
the ParlGov sample if comparable ideological coding can be achieved.

% ============================================================================
% NOVELTY REFERENCE SET (PLAN.TEX BACKBONE)
% ============================================================================

\section*{Novelty Reference Set}
\begin{itemize}[leftmargin=1.5em]
\item RD foundations and inference: \citep{lee2008randomized,calonico2014robust,cattaneo2015randomization}.
\item IV/LATE and sensitivity: \citep{imbens1994identification,conley2012plausibly}.
\item Core data sources: \citep{gwartney2025efw,clea2025,doring2021parlgov,dpi2020,vdem2025,vparty2024,worldbankwdi,pwt11}.
\end{itemize}

% ============================================================================
% REFERENCES
% ============================================================================

\newpage
\bibliographystyle{plainnat}
\bibliography{references}

% ============================================================================
% APPENDIX
% ============================================================================

\newpage
\appendix
\section{Appendix}

\subsection{Data Definitions}

\paragraph{Running Variable.}
Difference in vote share between the market bloc (Left-Right $\geq 5$) 
and the non-market bloc (Left-Right $< 5$).

\paragraph{EFW Index.}
The Fraser Institute Economic Freedom of the World index (0--10 scale). 
Area 5 refers to the Regulation component.

\paragraph{Cumulative Log GDP.}
$\sum_{s=1}^{h} (\log y_{t+s} - \log y_t)$, measuring the total log-point 
deviation from the election-year level.

\subsection{Country List}
The analysis includes elections from: Australia, Austria, Belgium, Canada, 
Denmark, Finland, France, Germany, Greece, Iceland, Ireland, Israel, 
Italy, Japan, Luxembourg, Netherlands, New Zealand, Norway, Portugal, 
Spain, Sweden, Switzerland, Turkey, United Kingdom.

\end{document}
